\documentclass[12pt]{article}
\usepackage[T1]{fontenc}
\usepackage[utf8]{inputenc}
\usepackage{lmodern}
\usepackage{textcomp}
\usepackage{enumitem}
\usepackage{geometry}
\geometry{margin=1in}
\begin{document}
\begin{center}
Answer Key - Business Administration - Graduate Level Assessment 3 \
\small (Answers are ordered exactly as in the question paper.)
\end{center}

\section*{Multiple Choice}
\begin{enumerate}[label=\arabic*.]
\item \textbf{Answer:} A
\\ \textit{Options:} A) \$925.00; B) \$890.42; C) \$650.00; D) \$785.53; E) \$600.00
\\ \textit{Note:} The correct answer of $925.00 is derived from applying the formula for compound interest, A = P(1 + r/n)^(nt), where P is the principal amount ($500), r is the annual interest rate (0.08), n is the number of times interest is compounded per year (4), and t is the number of years the money is invested (5), resulting in a total amount of \$925.00 after five years.
\item \textbf{Answer:} A
\\ \textit{Options:} A) 40 \%; B) 30\%; C) 20 \%; D) 80 \%; E) 45\%
\\ \textit{Note:} The percent markup based on the selling price is calculated by dividing the dollar markup ($80) by the selling price ($200), resulting in a markup of 40\%.
\item \textbf{Answer:} A
\\ \textit{Options:} A) 85.67; B) 90.12; C) 78.46; D) 82.19; E) 100.50
\\ \textit{Note:} The intrinsic value of the stock is calculated using the Gordon Growth Model, which determines the present value of future dividends growing at a constant rate, leading to an intrinsic value of $85.67 when a $5 dividend grows at 6\% and is discounted at a required return of 12\%.
\item \textbf{Answer:} D
\\ \textit{Options:} A) \$3035; B) \$2,985.75; C) \$3,000; D) \$3024.34; E) \$210
\\ \textit{Note:} The correct answer of \$3024.34 is derived by calculating the interest earned on the note for 60 days at 7\%, subtracting the discount at 5.5\% for the remaining period until maturity, resulting in the proceeds from the discounted note.
\item \textbf{Answer:} A
\\ \textit{Options:} A) \$1830.00; B) \$1785.00; C) \$1810.00; D) \$1799.50; E) \$1775.32
\\ \textit{Note:} The proceeds of the discounted note are calculated by first determining the interest for the remaining 70 days until maturity, which is $21, and then subtracting the bank's discount of $12 on the face value of $1800 at 6% for the same period, resulting in total proceeds of $1830.
\end{enumerate}

\section*{True / False}
\begin{enumerate}[label=\arabic*.]
\setcounter{enumi}{5}
\item \textbf{Answer:} True
\\ \textit{Options:} A) True; B) False
\\ \textit{Note:} The statement is true because a markup of 16(2/3)\% based on the selling price implies that the cost of the camera is 83.33\% of the selling price, resulting in a markup that aligns with the given percentage when calculated.
\item \textbf{Answer:} False
\\ \textit{Options:} A) True; B) False
\\ \textit{Note:} The term of discount for a note is typically calculated from the maturity date rather than the date of issue, which is why the statement is false.
\item \textbf{Answer:} False
\\ \textit{Options:} A) True; B) False
\\ \textit{Note:} The statement is false because the interest on a 60-day loan of $823 at an annual rate of 4% calculated using the exact time method is approximately $4.55, not \$4.96.
\item \textbf{Answer:} False
\\ \textit{Options:} A) True; B) False
\\ \textit{Note:} The statement is false because making deposits annually instead of monthly results in less frequent compounding, leading to a lower total amount accumulated over five years.
\item \textbf{Answer:} True
\\ \textit{Options:} A) True; B) False
\\ \textit{Note:} The statement is true because the interest calculated using the formula I = PRT (where I is interest, P is the principal amount, R is the rate, and T is the time in years) results in $7.21 for a principal of $677.21 at a 4\% annual interest rate over 90 days (or 0.25 years).
\end{enumerate}

\section*{Long Form Response}
\begin{enumerate}[label=\arabic*.]
\setcounter{enumi}{10}
\item \textbf{Answer:} To calculate the depreciation of his car during the seventh year using the sum of the years' digits (SYD) method, Bob Franklin would first determine the total useful life of the vehicle in years, which we will assume is 10 years for this example. The sum of the years' digits is calculated as 1 + 2 + 3 + ... + 10, resulting in a total of 55. For the seventh year, Bob would apply the fraction 4/55 (since 10-7+1 = 4) to the depreciable base of the car, which is the cost minus the salvage value. Assuming the depreciable base is $10,000, the depreciation expense for the seventh year would be $1,090.91. This method accelerates the depreciation in the earlier years, which can impact financial reporting by reducing taxable income sooner and providing a more conservative estimate of asset value on the balance sheet.
\\ \textit{Note:} correct because it accurately applies the sum of the years' digits method by calculating the total depreciation for the car's useful life, determining the appropriate fraction for the seventh year, and correctly identifying the depreciable base as the cost minus salvage value, thus demonstrating a thorough understanding of the method's application in financial reporting.
\item \textbf{Answer:} The underlying principles of the corporate governance Combined Code of Practice emphasize accountability, transparency, and fairness in the management of corporations. Integrity is essential in these governance frameworks as it fosters trust among stakeholders, including investors, employees, and customers. When corporate leaders demonstrate integrity, they create an environment where ethical behavior is prioritized, thereby enhancing the credibility of the organization and reducing the risk of fraudulent practices. Furthermore, integrity serves as a foundation for effective decision-making and risk management, ultimately contributing to sustainable business success and a positive corporate reputation. In this way, the principles of the Combined Code are not merely regulatory guidelines but essential components that underpin the ethical fabric of corporate governance.
\\ \textit{Note:} correct because it identifies the key principles of the Combined Code-accountability, transparency, and fairness-and explains how integrity underpins these principles by fostering trust and ethical behavior within the organization.
\end{enumerate}

\vfill
\noindent\textit{End of Answer Key}
\end{document}